% Source-reviewed research fragment, 8 October 2026.
% Not included in the authoritative manuscript or accepted Lean root.
% Requires amsmath, amssymb, amsthm and hyperref in a containing document.

\section{A rank-three entropy problem and two restricted bounds}
\label{sec:rank-three-entropy-research}

This note concerns a possible density extension of the finite atom description,
not a further proof of the obligatory classification.  Cycle and theta atoms
suggest asking whether the next trivalent core, $K_4$, admits an exact
edge-marginal entropy construction.  We obtain restricted statements and
identify two constructions that cannot establish the unrestricted result.
The arguments below have received analytical source review, not Lean
validation.  Their novelty and significance have not been established.

The initial equal-spoke, one-step compensation argument, the binary path
criterion and the obstruction examples were developed with a returned ChatGPT
Pro response.  The arbitrary-mismatch, unequal-spoke and two-diamond
refinements below were derived and checked separately during its audit.
The displayed model-version identity was unavailable.  These provenance
statements do not supply mathematical validation.

\subsection{The exact target and its boundary}

Replace each edge $ij$ of $K_4$ by a path of length $L_{ij}\geq1$,
with mutually disjoint interiors, and write $J=J(L)$.
The paths count edges, so
\[
 m=\sum_{ij}L_{ij},\qquad n=|V(J)|=m-2,\qquad m-n+1=3.
\]
We require $L_{ij}\equiv\varepsilon_i+\varepsilon_j\pmod2$ for some
branch colours $\varepsilon_i\in\{0,1\}$; thus $J$ is simple and bipartite.
Let $\eta$ be a symmetric exchangeable probability law on $\Omega^3$, where
$\Omega$ is finite.  Its pair and singleton marginals are $\lambda$ and
$\kappa$.  Restrict $\Omega$ to the positive support of $\kappa$; zero
entries of $\lambda$ remain allowed.  All logarithms are natural.
The question is whether there exists a probability law $\tau$ on the
\emph{actual} vertices of $J$ such that every oriented edge has marginal
$\lambda$ and
\begin{equation}
 D(\tau\Vert\kappa^{V(J)})\leq mD(\lambda\Vert\kappa^2).
 \label{eq:rank3-entropy-target}
\end{equation}
Every vertex is incident with an edge, so the prescribed edge laws fix all
singleton laws.  Hence the same question asks for
$H(\tau)\geq mH(\lambda)-(m+2)H(\kappa)$.
Mere feasibility is easy: use one sample from $\lambda$ on the two
bipartition classes.  It is the entropy budget that needs proof.

Known scalar density coverage must be removed before interpreting examples.  Sidorenko's
small-part theorem already covers $L=(2,3,1,1,1,2)$, which has parts of size
$4$ and $4$.\footnote{A.~Sidorenko, \emph{A correlation inequality for
bipartite graphs}, Graphs and Combinatorics 9 (1993), 201--204,
\url{https://doi.org/10.1007/BF02988307}.}
For $L_e=2a_e$, halving the paths gives a graph with $\sum_ea_e-2$
vertices and common kernel $T^2$.  Sidorenko's doubly-nonnegative theorem
therefore covers $\sum_ea_e\leq7$.\footnote{A.~Sidorenko,
\emph{Inequalities for doubly nonnegative functions}, Theorem 7.1 and
the subdivision discussion, \url{https://arxiv.org/abs/1905.08210}.}
Neither fact proves the target for arbitrary lengths and pair laws.

\subsection{A compensation bound with arbitrary arm mismatch}

Let $(\Omega,\mu)$ be a finite \emph{probability} space with
$\mu_x>0$ and $\sum_x\mu_x=1$.  Let $W$ be a symmetric kernel with
$0\leq W\leq1$ and constant row mean $p>0$:
\[
 \sum_y\mu_yW(x,y)=p.
\]
Kernel composition and inner products use $\mu$.  Put
$T=W/p$, $A=T^2$, $\Pi f=\int f\,d\mu$,
$K_j=A^j$, and $B_j=A^j-\Pi$ for $j\geq1$.
The self-adjoint Markov contraction $T$ fixes constants.  On the orthogonal
complement of constants, $A$ has eigenvalues $\theta_i\in[0,1]$.
For an orthonormal eigenbasis $\phi_0=1,\phi_1,\ldots$, write
\[
 B_j(x,y)=\sum_{i\geq1}\theta_i^j\phi_i(x)\phi_i(y),
 \qquad M_{ist}=\int\phi_i\phi_s\phi_t\,d\mu.
\]

For six positive integers $a_e$, the normalized density of the even
subdivision is
\[
 F(a)=\frac{t(J(2a),W)}{p^{2\sum_ea_e}}
     =\int_{\Omega^4}\prod_{e=vw\in E(K_4)}K_{a_e}(x_v,x_w)\,d\mu^4.
\]
Expand each factor as $1+B_{a_e}$.  Any term with a degree-one vertex
vanishes.  The remaining terms are the empty term, four triangles,
three four-cycles, six diamonds, and the full tetrahedron:
\begin{equation}
 F=1+\mathcal T+\mathcal C+\sum_eD_e+Q.
 \label{eq:rank3-complete-expansion}
\end{equation}
Here $D_e$ omits edge $e$.  Each triangle or cycle equals
$\sum_i\theta_i^{\sum a_e}\geq0$.
For example,
\[
 D_{cd}=\sum_{i,s,t\geq1}
 \theta_i^{a_{ab}}\theta_s^{a_{ac}+a_{bc}}
 \theta_t^{a_{ad}+a_{bd}}M_{ist}^2\geq0.
\]
Thus only $Q$ can be negative; its sign alone is not a density counterexample.

\begin{proposition}[A two-diamond compensation bound]
Suppose, in the order $ab,ac,ad,bc,bd,cd$,
\[
 a=(k,r+h,u,r,u,\ell),\qquad
 k\geq u\geq1,\quad r\geq\ell\geq1,\quad h\geq1.
\]
Set $d=r-\ell$ and
\[
 \gamma_{d,h}=\frac14\max_{0\leq x\leq1}
       \frac{x^d(1-x^h)^2}{1+x^h}.
\]
Here $x^0=1$, including at $x=0$.
If $p\geq\gamma_{d,h}$, then $t(J(2a),W)\geq p^m$, where
$m=2k+4r+4u+2h+2\ell$ and $|V(J(2a))|=m-2$.
The scale-invariant condition for bounded regular $W\geq0$ is
$p/\|W\|_\infty\geq\gamma_{d,h}$.
\end{proposition}

\begin{proof}
For fixed $z,w$, set
\[
 f(x)=B_{r+h}(x,z)B_u(x,w),\qquad
 g(x)=B_r(x,z)B_u(x,w),\qquad
 H=(f+g)/2,\quad\delta=(f-g)/2.
\]
Self-adjointness gives
\[
 Q+D_{cd}=\int K_\ell(z,w)\langle f,B_kg\rangle\,d\mu(z)d\mu(w)
 =\int K_\ell(z,w)
   \bigl(\langle H,B_kH\rangle-\langle\delta,B_k\delta\rangle\bigr)
   \,d\mu(z)d\mu(w).
\]
Since $K_\ell\geq0$ and $B_k\succeq0$, this is at least $-E$, where
$E=\int K_\ell\langle\delta,B_k\delta\rangle\,d\mu^2\geq0$.
Every positive power of $T$ is obtained by Markov averaging, so
$K_\ell=T^{2\ell}\leq1/p$ pointwise.  It follows that
$E\leq E_0/p$, with the exact expression
\[
 E_0=\frac14\sum_{i,s,t\geq1}
 \theta_i^k\theta_s^{2r}(1-\theta_s^h)^2
 \theta_t^{2u}M_{ist}^2.
\]
The two available diamonds are
\[
 D_{ac}+D_{bc}
 =\sum_{i,s,t\geq1}\theta_i^u\theta_s^{r+\ell}
       (1+\theta_s^h)\theta_t^{k+u}M_{ist}^2.
\]
For $x,y\in[0,1]$, $k\geq u$ implies
\[
 x^ky^{2u}+y^kx^{2u}
 \leq x^uy^{k+u}+y^ux^{k+u},
\]
because their difference, right minus left, is
$x^uy^u(x^u-y^u)(x^{k-u}-y^{k-u})\geq0$.
This includes $k=u$ and zero eigenvalues without division.
Symmetrize the sums in $i,t$, using $M_{ist}=M_{tsi}$.
The definition of $\gamma_{d,h}$ then gives
\[
 E_0\leq\gamma_{d,h}(D_{ac}+D_{bc}).
\]
Under the displayed density assumption, $E\leq D_{ac}+D_{bc}$.
Substitute in \eqref{eq:rank3-complete-expansion} and retain every other
nonnegative term.  This proves $F\geq1$.
For general bounded $W$, rescale by its positive supremum.
\end{proof}

The threshold is explicit.  For $d=0$, $\gamma_{0,h}=1/4$.
For $d>0$, put $b=d/h$ and
\[
 y_* =\frac{2b}{3+\sqrt{9+4b(b+1)}}.
\]
Then
\[
 \gamma_{d,h}=\frac14\,
       \frac{y_*^b(1-y_*)^2}{1+y_*}.
\]
Indeed, differentiating $y^b(1-y)^2/(1+y)$ gives the unique interior
stationary point satisfying $b-3y-(b+1)y^2=0$; the endpoint values vanish.
A simpler, weaker threshold is
\[
 c_{d,h}=\frac14\max_{0\leq x\leq1}x^d(1-x^h)^2
 =\begin{cases}
 1/4,&d=0,\\
 \displaystyle\frac{h^2}{(d+2h)^2}
       \left(\frac{d}{d+2h}\right)^{d/h},&d>0.
 \end{cases}
\]
It specializes, for $h=1$, to $c_{d,1}=d^d/(d+2)^{d+2}$.
For example, lengths $(2,6,2,4,2,2)$ are covered at $p\geq1/27$;
the two-diamond threshold is smaller.  These remain \emph{restricted}
weighted-host inequalities.  At any fixed tuple, sufficiently sparse hosts
remain untreated, so the proposition does not establish
\eqref{eq:rank3-entropy-target} through a universal type-host argument.

\subsection{An exact binary entropy construction}

Let $\kappa$ be uniform on $\{-1,1\}$ and
$\lambda_\rho(s,t)=(1+\rho st)/4$.  Write
\[
 I(\rho)=\frac{1+\rho}{2}\log(1+\rho)
          +\frac{1-\rho}{2}\log(1-\rho).
\]
The following criterion concerns prescribed marginals, not merely a
homomorphism count with one common edge potential.

\begin{proposition}[A binary path criterion]
Let $G$ be a finite simple graph and $0\leq r<1$.
For each edge $e=vw$, let $\mathcal P_e$ be the simple $v$--$w$ paths in $G-e$.
If $\sum_{\pi\in\mathcal P_e}r^{|\pi|}\leq r$ for every $e$, there exists
a law $\tau$ on $\{-1,1\}^{V(G)}$ with every edge marginal $\lambda_r$ and
$D(\tau\Vert\kappa^{V(G)})\leq |E(G)|I(r)$.
For bipartite $G$, the same holds with $\lambda_{-r}$.
\end{proposition}

\begin{proof}
For $r=0$, use independent spins.  Otherwise put $b=\operatorname{atanh}r$
and minimize the continuous function
\[
 \Phi(\beta)=\log Z(\beta)-r\sum_e\beta_e,\qquad
 Z(\beta)=2^{-|V(G)|}\sum_\sigma
      \exp\left(\sum_{e=vw}\beta_e\sigma_v\sigma_w\right)
\]
on the compact box $[0,b]^{E(G)}$.
With $t_e=\tanh\beta_e$, the finite high-temperature expansion is
\[
 Z(\beta)=\prod_e\cosh\beta_e
       \sum_{A\subseteq E(G):\,\partial A=\varnothing}\prod_{e\in A}t_e
       \geq\prod_e\cosh\beta_e.
\]
In $G-e$, its endpoint correlation $c_e$ lies between zero and
$\sum_{\pi\in\mathcal P_e}\prod_{f\in\pi}t_f\leq r$.
To justify the path bound, every set with odd boundary $\{v,w\}$ contains
a simple $v$--$w$ path.  Remove one such path; the remaining set is even.
Summing over all paths and even remainders overcounts the nonnegative
odd-boundary numerator and gives the asserted bound after division by
the even-subgraph denominator.
Reinserting $e$ gives correlation $(t_e+c_e)/(1+t_ec_e)$ and
$\partial_e\Phi=\mathbb E(\sigma_v\sigma_w)-r$.
At a lower box face this derivative is nonpositive, whereas a constrained
minimum requires it to be nonnegative.  At an upper face it is
nonnegative, whereas a minimum requires it to be nonpositive.
All derivatives thus vanish at a minimizer $\beta^*$.
The resulting Gibbs law has correlation $r$ on each edge.
Global spin-flip symmetry fixes uniform singletons, hence the exact law
$\lambda_r$ on every edge.  Finally,
\[
 D(\tau\Vert\kappa^{V(G)})
 =r\sum_e\beta_e^*-\log Z(\beta^*)
 \leq\sum_e\bigl(r\beta_e^*-\log\cosh\beta_e^*\bigr)
 \leq |E(G)|I(r).
\]
The last bound is the elementary Legendre transform of $\log\cosh$.
Flip the spins on one bipartition class to obtain $-r$.
\end{proof}

In a $K_4$ subdivision, deletion of an actual edge leaves exactly four
alternate simple paths between its endpoints: after the forced path tails,
they are the core routes $acb,adb,acdb,adcb$ for the relevant branch edge $ab$.
For a simple bipartite graph each has length at least three.
Thus $4r^3\leq r$ proves the criterion for $r\leq1/2$.
The uniform binary pair law has an exchangeable triple extension exactly
when $-1/3\leq\rho\leq1$: necessity follows from
$(\sigma_1+\sigma_2+\sigma_3)^2\geq1$, and sufficiency from
\[
 \eta_\rho(x,y,z)=\frac18\bigl(1+\rho(xy+xz+yz)\bigr).
\]
Consequently \eqref{eq:rank3-entropy-target} holds for this specified
binary family when $-1/3\leq\rho\leq1/2$.
This source-reviewed result is not being claimed as novel; exact
literature reconciliation is a separate obligation.
The argument belongs to established Ising and Bethe methodology.
Ruozzi proves a partition-function lower bound for binary
log-supermodular models,\footnote{N.~Ruozzi,
\emph{The Bethe Partition Function of Log-supermodular Graphical Models},
Advances in Neural Information Processing Systems 25 (2012),
\url{https://proceedings.neurips.cc/paper_files/paper/2012/file/03afdbd66e7929b125f8597834fa83a4-Paper.pdf}.}
whereas attractive Ising fitting permits excess fitted correlation on
zero-coupling edges.\footnote{S.~Lauritzen, C.~Uhler and P.~Zwiernik,
\emph{Total positivity in exponential families with application to binary variables},
Corollary 3.7 and Remark 3.9,
\url{https://web.math.ku.dk/~lauritzen/papers/AOS2007.pdf}.}
Neither statement alone supplies all the exact prescribed marginals here;
the path criterion discharges that additional obligation.

\subsection{Two limits on proposed constructions}

\paragraph{Direct mixture inheritance fails.}
Mix, with weights $1/2,1/2$, the iid uniform binary triple and the triple
that is $000$ or $111$ with equal probabilities.  The result is strictly
positive, with masses $5/16$ at $000,111$ and $1/16$ elsewhere, and pair
law $\lambda_{1/2}$.
For the iid component a budget-satisfying extension is forced to be iid,
since its allowed divergence is zero.  For the diagonal component,
connectedness forces the unique law supported on $0^n,1^n$.
Their direct mixture is therefore
\[
 \tau_{\rm mix}=\tfrac12\kappa^n
       +\tfrac14\delta_{0^n}+\tfrac14\delta_{1^n},\qquad
 D(\tau_{\rm mix}\Vert\kappa^n)
       =(\tfrac12+2^{-n})\log(2^{n-1}+1)-\log2.
\]
For $J_0=J(2,4,3,2,3,1)$, $n=13,m=15$, and
\[
 D(\tau_{\rm mix}\Vert\kappa^{13})>5\log2
       >15I(1/2)=15(\tfrac34\log3-\log2).
\]
The first comparison uses $4097>2^{12}$ and $4097/8192>1/2$;
the second is equivalent to $2^{80}>3^{45}$, which follows already
from $2^{16}>3^9$.
The binary construction above supplies a \emph{different} successful
extension for this same law.  Thus this is a counterexample to direct
mixture inheritance, not to the target.

For components sharing $\kappa$, let
$\delta_z=mD(\lambda_z\Vert\kappa^2)-D(\tau_z\Vert\kappa^n)$.
The exact chain rule is
\[
 D(\tau_{\rm mix}\Vert\kappa^n)-mD(\lambda_{\rm mix}\Vert\kappa^2)
 =mI(Z;X_v,X_w)-I(Z;X_V)-\sum_z\alpha_z\delta_z.
\]
The example violates the desired sign even after component slack is included.
It does not decide convexity of the class of pair laws admitting some
successful extension.

\paragraph{Commuting positive kernels are not common powers.}
Index six signs $r_e\in\{-1,1\}$ by the edges of $K_4$, and put
$S_v(r)=\prod_{e\ni v}r_e$.
On these sign vectors define
\[
 \mu(r)=2^{-6}\left(1+\frac{S_a(r)+S_b(r)+S_c(r)-S_d(r)}2\right).
\]
Since $S_aS_bS_cS_d=1$, the signed sum is $2$ or $-2$.
After removing null states this is uniform on $32$ states.
The six coordinates are centered and orthonormal: averaging their
degree-one and degree-two characters against this density gives zero.
Each $K_e(x,y)=1+r_e(x)r_e(y)$ is a symmetric nonnegative Markov
projection; the six projections commute.
In the coloured $K_4$ density every proper nonempty centered term
vanishes at a vertex of degree one or two.  The full term is
\[
 (\mathbb E S_a)(\mathbb E S_b)(\mathbb E S_c)(\mathbb E S_d)
 =(\tfrac12)^3(-\tfrac12)=-\tfrac1{16}.
\]
The density is $15/16<1$.
Positive powers of one self-adjoint operator have a common nullspace,
whereas these projections have distinct nullspaces.
Each individual pair law even has an exchangeable triple extension:
sample its common coordinate sign and then three iid states conditional
on that sign.  What is absent is the \emph{same} pair law and common-power
relationship.  This exact example forbids that relaxation, not the target.

\subsection{The unresolved common-power comparison}

For $a=(k,r+h,u,r,u,\ell)$, the sufficient inequality
\[
 \int K_\ell(z,w)\langle
 B_{r+h}(\cdot,z)B_u(\cdot,w),
 B_k[B_r(\cdot,z)B_u(\cdot,w)]\rangle\,d\mu^2\geq0
\]
would settle this particular all-even family on every normalized weighted
host: its left side is $Q+D_{cd}$.
It is unproved here and stronger than what the compensation proposition
needs.  The factor $K_\ell(z,w)$ changes the inner-product measure;
positivity and self-adjointness in the original measure do not justify
a change of measure.  The following example shows that this is an actual
obstruction to the arbitrary-vector argument, even for common powers.

\paragraph{Weighted positivity fails on an admissible common-power host.}
Take three states with uniform measure, $p=6/13$, and
\[
 W=\begin{pmatrix}
 1&1/13&4/13\\
 1/13&1&4/13\\
 4/13&4/13&10/13
 \end{pmatrix},\qquad
 P=\begin{pmatrix}
 13/18&1/18&2/9\\
 1/18&13/18&2/9\\
 2/9&2/9&5/9
 \end{pmatrix}.
\]
Here $W$ is symmetric, strictly positive, bounded by one, and has row
mean $p$.  The matrix $P$ represents the operator $T=W/p$;
the kernel is $3P$, because composition uses the uniform probability
measure.  Thus $A$ is represented by $P^2$, with eigenvalues
$1,4/9,1/9$.  Set $v=(-1,-3,4)$ and choose the first state as $w$.
Since $v$ is mean zero and both nonconstant eigenvalues are positive,
$v$ belongs to the range of $B_1$.
Direct rational arithmetic gives
\[
 A^3v=\frac1{729}(62,-66,4),\qquad
 K_1(\cdot,w)=\left(\frac{31}{18},\frac7{18},\frac89\right),
\]
and hence
\[
 \int K_1(z,w)v(z)(A^3v)(z)\,d\mu(z)
       =-\frac{140}{19683}<0.
\]
Common powers and smoothing therefore do not imply weighted positivity
for arbitrary mean-zero inputs.  This does not refute $Q+D_{cd}\geq0$:
the inputs in that expression are products of eigenfunctions and kernel
sections, summed with prescribed spectral weights.  Nor does it refute
$F\geq1$, which retains the other nonnegative terms in
\eqref{eq:rank3-complete-expansion}.

\paragraph{A comparison retaining all proper terms.}
For each nonconstant eigenfunction define
\[
 R_s=\int\phi_s(a)\phi_s(c)
 K_k(a,b)K_u(a,d)K_r(b,c)K_u(b,d)K_\ell(c,d)\,d\mu^4.
\]
Expanding only the $ac$ kernel gives the exact identity
\[
 F=R_0+\sum_{s\geq1}\theta_s^{r+h}R_s,
\]
where $R_0$ denotes the same integral with $\phi_0=1$.
The remaining five-edge graph is an uncentered diamond, so
\[
 R_0=1+\sum_{i\geq1}\theta_i^{k+2u}
       +\sum_{i\geq1}\theta_i^{r+u+\ell}
       +\sum_{i\geq1}\theta_i^{k+r+u+\ell}+D_{ac}\geq1.
\]
Group the distinct positive nonconstant eigenvalues as
$\eta_1>\cdots>\eta_q>0$, and put
$c_j=\eta_j^r\sum_{\theta_s=\eta_j}R_s$ and
$S_j=\sum_{v=1}^jc_v$.  These grouped sums do not depend on the choice
of basis inside a repeated eigenspace.  For $q\geq1$, Abel summation gives
\[
 F-R_0=\sum_{j=1}^{q-1}(\eta_j^h-\eta_{j+1}^h)S_j
             +\eta_q^hS_q.
\]
If $q=0$, every positive kernel power is $\Pi$ and $F=R_0=1$.
Thus positivity of every spectral prefix $S_j$ would suffice.  That
positivity is not proved here; even its failure would invalidate only
this sufficient route, since $R_0-1$ can compensate negative coefficients.
The necessary comparison is
$\sum_{s\geq1}\theta_s^{r+h}R_s\geq-(R_0-1)$, with the actual common
host and moment constraints retained.

Six independent exponents and odd-parity paths
require still further control.  The unrestricted entropy theorem,
formalization of these analytical results, global novelty and B-level
significance remain open obligations rather than conclusions of this note.
